\chapter{Desarrollo}

\section{Diseño general del sistema}

El desarrollo se realizó sobre una versión anterior ya funcional del proyecto. Este consiste en un sistema de gestión de deportistas, que permite cargar datos desde un archivo CSV, ordenarlos según distintos criterios y mostrar resultados específicos como el ranking de los mejores deportistas. Este se basó en la idea de evaluar el rendimiento de diferentes algoritmos de ordenamiento y búsqueda secuenciales, los cuales se mantuvieron en la versión actual.

\subsection{Deportistas y generación de datos}
Los deportistas se representan mediante una estructura que contiene: \textit{ID}, \textit{nombre}, \textit{puntaje} y \textit{competencias}. La generación de datos se realiza mediante la opción \textit{-g <numero>} del programa, que crea un archivo CSV con el número especificado de deportistas. Cada deportista tiene un ID único, un nombre generado aleatoriamente, y un puntaje y cantidad de competencias también aleatorios dentro de rangos predefinidos. Esto permite probar los algoritmos con conjuntos de datos de diferentes tamaños y características.

\subsection{Funcionalidades}
El programa implementa una interfaz de línea de comandos, que ofrece las siguientes funcionalidades principales:
\begin{itemize}
    \item \textbf{Generación de datos:} Permite crear archivos CSV con un número específico de deportistas aleatorios.
    \item \textbf{Ordenamiento:} Dispone de un menú en el que el usuario puede ordenar los deportistas, seleccionando el algoritmo, el criterio de ordenamiento (ID, puntaje, competencias) y el orden (ascendente o descendente).
    \item \textbf{Búsqueda:} Permite buscar deportistas por ID, seleccionando el algoritmo de búsqueda. Si se selecciona la búsqueda binaria por rango, sin embargo este buscará por puntaje, mostrando todos los deportistas que tengan el puntaje ingresado.
    \item \textbf{Ranking:} Permite mostrar el ranking de los mejores \(N\) deportistas según puntaje, utilizando Quick Select para obtener los resultados de manera eficiente.
    \item \textbf{Búsqueda de k-ésimo elemento:} Permite obtener el deportista ubicado en la posición \(k\) según un criterio específico, también utilizando Quick Select.
    \item \textbf{Búsqueda por rango de puntajes:} Permite mostrar todos los deportistas cuyo puntaje se encuentre dentro de un intervalo definido por el usuario, utilizando búsqueda binaria por rango.
    \item \textbf{Benchmarks:} En el programa se incluyen 4 benchmarks de experimentación, los cuales miden el tiempo de ejecución de los algoritmos de ordenamiento, búsqueda y selección sobre conjuntos de datos de diferentes tamaños y características, guardando los resultados en archivos CSV para su posterior análisis.
\end{itemize}

\section{Algoritmos de Ordenamiento}

\subsection{Merge Sort}

La implementación de Merge Sort se realizó mediante una estrategia recursiva. El algoritmo divide sucesivamente el arreglo de deportistas en dos mitades hasta alcanzar el caso base, correspondiente a subarreglos de un único elemento, los cuales ya se consideran ordenados.

La función recursiva calcula el punto medio utilizando la expresión:

\[
\texttt{mid = left + (right - left) / 2}
\]

evitando posibles errores cuando los índices son muy grandes que podrían ocurrir con la expresión clásica:

\[
\texttt{(left + right) / 2}
\]

Una vez ordenadas ambas mitades, se utiliza una función auxiliar \texttt{merge} encargada de fusionarlas en un único segmento ordenado. Para ello, se reservan dinámicamente dos arreglos temporales (\texttt{L} y \texttt{R}) que almacenan copias de las sublistas izquierda y derecha respectivamente.

Durante la mezcla, los elementos son comparados mediante la función \texttt{compare\_by\_criteria}, permitiendo ordenar según distintos criterios del sistema (ID, puntaje o cantidad de competencias). Además, la función \texttt{should\_precede} abstrae la lógica de comparación para soportar tanto orden ascendente como descendente.

El proceso de mezcla recorre simultáneamente ambos arreglos temporales e inserta en el arreglo original el elemento que debe preceder al otro según el criterio seleccionado. Cuando una de las mitades se agota, los elementos restantes de la otra mitad son copiados directamente.

Debido al uso de memoria auxiliar para las sublistas temporales, el algoritmo requiere una complejidad espacial adicional de:

\[
O(n)
\]

mientras que su complejidad temporal se mantiene estable en:

\[
O(n \log n)
\]

tanto en el caso promedio como en el peor caso.

\subsection{Merge Sort Optimizado}

La versión optimizada de Merge Sort conserva la misma estructura recursiva y el mismo procedimiento de mezcla que la implementación estándar, pero incorpora una optimización híbrida basada en Insertion Sort para mejorar el rendimiento práctico sobre subarreglos pequeños.

Durante la recursión, antes de seguir dividiendo el arreglo, se evalúa el tamaño del subarreglo actual mediante la condición:

\[
\texttt{right - left + 1 <= MERGE\_SORT\_THRESHOLD}
\]

Si el tamaño es menor o igual al umbral definido, el algoritmo deja de subdividir y utiliza \texttt{Insertion Sort} sobre ese segmento.

Esta optimización se basa en que Insertion Sort posee un costo constante muy bajo y un excelente desempeño en arreglos pequeños o parcialmente ordenados, evitando el overhead asociado a múltiples llamadas recursivas y asignaciones dinámicas de memoria.

La implementación de \texttt{Insertion Sort} desplaza los elementos mayores hacia la derecha hasta encontrar la posición correcta del elemento actual (\texttt{key}). Al igual que en Merge Sort, las comparaciones se realizan mediante la función \texttt{should\_precede}, manteniendo compatibilidad con múltiples criterios y órdenes de ordenamiento.

La operación de mezcla sigue utilizando arreglos auxiliares dinámicos para fusionar las mitades ya ordenadas. Además, la implementación incluye validaciones de memoria después de cada llamada a \texttt{malloc}. En caso de error de asignación, se liberan inmediatamente los recursos reservados mediante \texttt{free}, evitando fugas de memoria.

Aunque la complejidad asintótica permanece en:

\[
O(n \log n)
\]

la reducción de llamadas recursivas y operaciones de memoria sobre subarreglos pequeños produce una mejora significativa en el rendimiento real del algoritmo.

\subsection{Quick sort}

\subsubsection{Partición de Lomuto}
La partición es la operación central del algoritmo. Recibe el subarreglo delimitado por low y high, tomando siempre el último elemento como pivote. Mantiene un índice i que comienza antes del inicio del subarreglo y marca el límite de los elementos que ya se sabe que deben ir a la izquierda del pivote. El índice j recorre todos los elementos del subarreglo excepto el pivote. Cada vez que encuentra un elemento que pertenece a la izquierda, avanza i e intercambia ese elemento con deportistas[i]. Al terminar el recorrido, coloca el pivote en i + 1, que es su posición definitiva, y retorna ese índice.

\subsubsection{Variantes de selección de pivote}
Todas las variantes comparten exactamente el mismo código de partición y recursión. La única diferencia entre ellas es qué hacen antes de partir, y todas terminan dejando el pivote elegido en la última posición para que Lomuto pueda operar siempre de la misma manera.

\subsubsection{Variante de último elemento}
La variante de último elemento no realiza ninguna acción previa, ya que el pivote ya está donde Lomuto lo espera. Es la más simple pero la más vulnerable: en arreglos ya ordenados el pivote siempre cae en un extremo, lo que degenera el rendimiento a $O(n^2)$.

\subsubsection{Variante de primer elemento}
La variante de primer elemento simplemente intercambia el primer elemento con el último antes de partir. Resuelve el caso del arreglo invertido pero sigue siendo vulnerable al arreglo ya ordenado, que es el peor caso contrario.

\subsubsection{Variante de elemento aleatorio}
La variante de elemento aleatorio elige un índice al azar dentro del rango y lo mueve al final. Al no depender de la disposición actual del arreglo, rompe los patrones adversos que afectan a los pivotes fijos y garantiza un buen rendimiento promedio sin importar el orden inicial de los datos.

\subsubsection{Variante de mediana de tres}
La variante de mediana de tres examina el primer, el del medio y el último elemento, los ordena entre sí mediante comparaciones directas, y usa la mediana como pivote moviéndola al final. Es la variante más robusta: evita los peores casos de los pivotes fijos, produce particiones más balanceadas que el pivote aleatorio en la mayoría de los casos, y no depende de un generador de números aleatorios.

\subsubsection{Estructura recursiva}
Cada función sigue la misma estructura: si \( low < high \), se posiciona el pivote, se particiona el subarreglo obteniendo el índice final del pivote \( pi \), y se llama recursivamente sobre los intervalos \([low, pi-1]\) y \([pi+1, high]\). El caso base es implícito: cuando \( low \geq high \), el subarreglo tiene un solo elemento o está vacío, por lo que ya está ordenado y no hay nada que hacer.

La complejidad en el caso promedio es:

\[
O(n \log n)
\]

El peor caso es:

\[
O(n^2)
\]

y ocurre cuando el pivote siempre queda en un extremo, situación que las variantes de pivote aleatorio y mediana de tres mitigan considerablemente.


\section{Algoritmos de Búsqueda y Selección}

\subsection{Búsqueda Binaria Recursiva}

La búsqueda binaria recursiva se implementa mediante la función \texttt{recursive\_binary\_search}. La función recibe cinco parámetros:

\begin{itemize}
    \item \textbf{deportistas:} Puntero al arreglo de estructuras deportista, que debe estar ordenado por ID\@.
    \item \textbf{left:} Índice del límite izquierdo del rango actual de búsqueda\@.
    \item \textbf{right:} Índice del límite derecho del rango actual de búsqueda\@.
    \item \textbf{targetId:} El ID del deportista que se desea encontrar\@.
\end{itemize}

La implementación se divide en tres partes principales, siguiendo la estrategia de divide y vencerás:

\subsubsection{Caso base}
El primer condicional evalúa si \texttt{left > right}. Esta condición indica que el rango de búsqueda se ha agotado sin encontrar el elemento, por lo que la función retorna \texttt{-1} para señalar que el ID no existe en el arreglo.

\subsubsection{División}
Si el rango aún contiene elementos (\texttt{left} $\leq$ \texttt{right}), el algoritmo calcula el punto medio del rango actual:

\[
\texttt{mid = left + (right - left) / 2}
\]

Esta expresión evita desbordamientos que podrían ocurrir con la forma clásica \texttt{(left + right) / 2} en arreglos muy grandes. El índice \texttt{mid} divide el problema en dos subproblemas de aproximadamente igual tamaño.

\subsubsection{Conquista}
Una vez dividido el problema, se compara el elemento en la posición \texttt{mid} con el \texttt{targetId}:

\begin{itemize}
    \item \textbf{Elemento encontrado:} Si \texttt{deportistas[mid]->id == targetId}, se retorna el índice \texttt{mid} inmediatamente.
    
    \item \textbf{Buscar en la mitad izquierda:} Si \texttt{deportistas[mid]->id > targetId}, el elemento buscado debe estar en la mitad izquierda, ya que el arreglo está ordenado. Se realiza una llamada recursiva:
    \[
    \texttt{recursive\_binary\_search(deportistas, left, mid - 1, criteria, targetId)}
    \]
    
    \item \textbf{Buscar en la mitad derecha:} Si \texttt{deportistas[mid]->id < targetId}, el elemento está en la mitad derecha. Se realiza la llamada recursiva:
    \[
    \texttt{recursive\_binary\_search(deportistas, mid + 1, right, criteria, targetId)}
    \]
\end{itemize}

\subsection{Búsqueda exponencial}

La búsqueda exponencial implementa la estrategia divide y vencerás combinando dos fases: expansión exponencial para localizar un rango de búsqueda y luego búsqueda binaria recursiva dentro de ese rango.

\subsubsection{Fase de expansión exponencial}

El algoritmo comienza verificando si el elemento objetivo se encuentra en la primera posición. Si es así, retorna el índice 0 inmediatamente, logrando el mejor caso \(O(1)\).

En caso contrario, inicia la fase de expansión con un límite \texttt{bound} establecido en 1. En cada iteración, duplica el valor del límite:

\[
\texttt{bound = bound * 2}
\]

Durante cada iteración, verifica si el elemento en la posición \texttt{bound} es mayor u igual al objetivo. Si lo es, ha encontrado un rango que contiene potencialmente al elemento buscado. El rango identificado se extiende desde \(\texttt{bound/2}\) hasta \texttt{bound} (o hasta \texttt{length - 1} si \texttt{bound} excede la longitud del arreglo).

La cantidad de iteraciones en esta fase es proporcional a \(\log i\), donde \(i\) es la posición del elemento buscado, ya que cada duplicación del límite acota exponencialmente el espacio de búsqueda.

\subsubsection{Fase de búsqueda binaria recursiva}

Una vez identificado el rango, el algoritmo aplica búsqueda binaria de forma \textbf{completamente recursiva} sobre el intervalo \([\texttt{bound/2}, \texttt{bound}]\). La implementación utiliza una función auxiliar \texttt{exponencial\_binary\_search\_recursive} que implementa divide y vencerás como se describió anteriormente.

\subsection{Búsqueda por Interpolación}

La búsqueda por interpolación se implementó como una alternativa optimizada a la búsqueda binaria para arreglos numéricos ordenados, utilizando una estrategia de divide y vencerás mediante \textbf{división inteligente}.

A diferencia de la búsqueda binaria, que siempre divide el rango por la mitad, este algoritmo estima una posición probable del elemento objetivo basándose en los valores reales en los extremos del rango, permitiendo una división más cercana al objetivo.

\subsubsection{Caso base}
Si \texttt{low > high}, el rango se ha agotado sin encontrar el elemento, por lo que retorna \texttt{-1}.

\subsubsection{División}
En lugar de calcular \texttt{mid = (left + right) / 2} como en búsqueda binaria, la interpolación estima la posición usando la fórmula:

\begin{equation}
    pos = low + \left[ \frac{high - low}{highValue - lowValue} \cdot (targetValue - lowValue) \right]
\end{equation}

Donde:
\begin{itemize}
    \item \textbf{low / high:} Son los índices de los límites del rango actual.
    \item \textbf{lowValue / highValue:} Son los valores en dichos índices según el criterio de búsqueda.
    \item \textbf{targetValue:} Es el valor que se busca.
\end{itemize}

Esta fórmula mapea linealmente el valor buscado dentro del rango conocido de valores, produciendo una estimación más cercana al objetivo en datos con distribución uniforme. Esta es la división del problema: en lugar de un corte aritmético fijo, es un corte proporcional al valor.

\subsubsection{Conquista}
Una vez estimada la posición, se compara el valor encontrado con el objetivo:

\begin{itemize}
    \item \textbf{Elemento encontrado:} Si los valores coinciden, retorna el índice.
    
    \item \textbf{Buscar en la mitad derecha:} Si el valor en \texttt{pos} es menor que el objetivo, se recursiona:
    \[
    \texttt{interpolation\_search\_recursive(deportistas, pos + 1, high, criteria, targetValue)}
    \]
    
    \item \textbf{Buscar en la mitad izquierda:} Si el valor en \texttt{pos} es mayor que el objetivo, se recursiona:
    \[
    \texttt{interpolation\_search\_recursive(deportistas, low, pos - 1, criteria, targetValue)}
    \]
\end{itemize}

\subsection{Búsqueda por Rango}

La búsqueda por rango implementa divide y vencerás para encontrar \textbf{todos} los elementos que coinciden con un valor específico, retornando los índices del primer y último elemento que poseen ese valor. Esto es especialmente útil cuando múltiples deportistas comparten el mismo puntaje o criterio de búsqueda.

\subsubsection{Estructura recursiva}

El algoritmo utiliza tres funciones recursivas auxiliares que implementan divide y vencerás en cada fase:

\textbf{1. Fase de localización inicial (find\_any\_recursive):}

Realiza una búsqueda binaria recursiva estándar para encontrar \textbf{cualquier} elemento que coincida con el valor objetivo. Una vez que encuentra un elemento coincidente, retorna su índice:

\begin{itemize}
    \item \textbf{Caso base:} Si \texttt{low > high}, no se encontró nada, retorna $-1$.
    \item \textbf{División:} Calcula \(\texttt{mid = low + (high - low) / 2}\).
    \item \textbf{Conquista:} 
    \begin{itemize}
        \item Si \texttt{deportistas[mid]} coincide con el objetivo, retorna \texttt{mid}.
        \item Si el valor en \texttt{mid} es menor, recursa en la mitad derecha.
        \item Si es mayor, recursa en la mitad izquierda.
    \end{itemize}
\end{itemize}

\textbf{2. Fase de expansión izquierda (find\_leftmost\_recursive):}

Una vez localizado cualquier elemento coincidente, esta función busca recursivamente el \textbf{primer} elemento con ese valor expandiendo hacia la izquierda:

\begin{itemize}
    \item \textbf{Caso base:} Si \texttt{low > high}, retorna $-1$.
    \item \textbf{División:} Calcula \(\texttt{mid = low + (high - low) / 2}\).
    \item \textbf{Conquista:} 
    \begin{itemize}
        \item Si \texttt{deportistas[mid]} coincide, continúa buscando en la mitad izquierda pero \textbf{recuerda este índice} como el candidato más a la izquierda.
        \item Si encuentra un elemento más a la izquierda, lo retorna; de lo contrario, retorna el índice actual.
        \item Si el valor es menor, busca a la derecha.
        \item Si es mayor, busca a la izquierda.
    \end{itemize}
\end{itemize}

\textbf{3. Fase de expansión derecha (find\_rightmost\_recursive):}

Busca recursivamente el \textbf{último} elemento con el valor objetivo expandiendo hacia la derecha:

\begin{itemize}
    \item \textbf{Caso base:} Si \texttt{low > high}, retorna $-1$.
    \item \textbf{División:} Calcula \(\texttt{mid = low + (high - low) / 2}\).
    \item \textbf{Conquista:} 
    \begin{itemize}
        \item Si \texttt{deportistas[mid]} coincide, continúa buscando en la mitad derecha pero \textbf{recuerda este índice} como el candidato más a la derecha.
        \item Si encuentra un elemento más a la derecha, lo retorna; de lo contrario, retorna el índice actual.
        \item Si el valor es menor, busca a la derecha.
        \item Si es mayor, busca a la izquierda.
    \end{itemize}
\end{itemize}

\subsubsection{Tipo de dato retornado}

La función retorna una estructura \texttt{Range} que contiene dos campos:

\begin{itemize}
    \item \texttt{start:} Índice del primer elemento que coincide con el valor objetivo.
    \item \texttt{end:} Índice del último elemento que coincide con el valor objetivo.
\end{itemize}

Si no se encuentra ningún elemento, ambos campos se establecen en $-1$.

\subsection{Quick Select}

Para obtener el elemento ubicado en una posición específica sin ordenar completamente el arreglo, se implementó el algoritmo \textit{Quick Select} como una versión recursiva de divide y vencerás. Este algoritmo reutiliza la lógica de partición de Quick Sort, pero únicamente continúa la búsqueda sobre la partición que contiene el elemento deseado, reduciendo significativamente el trabajo realizado.

\subsubsection{Selección de pivote con mediana de tres}

Antes de particionar, el algoritmo utiliza la estrategia de \textbf{mediana de tres} para elegir el pivote:

\begin{itemize}
    \item Se examinan tres elementos: el primero (\texttt{left}), el del medio, y el último (\texttt{right}).
    \item Se ordenan estos tres elementos entre sí mediante comparaciones directas.
    \item La mediana de estos tres se coloca en la última posición (\texttt{right}), donde la partición la espera.
\end{itemize}

Esta estrategia evita los peores casos de pivotes fijos (primer o último elemento) y produce particiones más balanceadas que un pivote aleatorio en la mayoría de los casos prácticos.

\subsubsection{Partición de Lomuto}

La operación principal es la función \texttt{partition}, basada en el esquema de Lomuto. Durante el recorrido del subarreglo:

\begin{itemize}
    \item Se mantiene un índice \(i\) que marca el límite entre elementos que deben ir a la izquierda del pivote y los que no.
    \item Se comparan todos los elementos (excepto el pivote) con el pivote, decidiendo si preceden o no según el criterio y orden seleccionados.
    \item Los elementos que deben estar a la izquierda se intercambian hacia esa posición.
    \item Al final, el pivote se coloca en su posición definitiva.
\end{itemize}

La comparación se realiza mediante \texttt{compare\_by\_criteria}, permitiendo trabajar con múltiples criterios (ID, puntaje, competencias), y el orden (ascendente o descendente) se controla mediante el parámetro \texttt{order}.

\subsubsection{Estructura recursiva de divide y vencerás}

La función \texttt{quick\_select\_recursive} implementa el ciclo de divide y vencerás:

\begin{itemize}
    \item \textbf{Caso base:} Si el subarreglo contiene un único elemento (\texttt{left == right}), se retorna directamente sin más procesamiento.
    
    \item \textbf{División:} Se aplica mediana de tres y se particiona el subarreglo, obteniendo el índice del pivote (\texttt{pivotIndex}).
    
    \item \textbf{Conquista:} Se evalúan tres casos:
    \begin{itemize}
        \item Si \(k =\) \texttt{pivotIndex}, el elemento buscado fue encontrado en su posición correcta.
        \item Si \(k <\) \texttt{pivotIndex}, el elemento está en la mitad izquierda, se recursa:
        \[
        	texttt{quick\_select\_recursive(arr, left, pivotIndex - 1, k, \ldots)}
        \]
        \item Si \(k >\) \texttt{pivotIndex}, el elemento está en la mitad derecha, se recursa:
        \[
        	texttt{quick\_select\_recursive(arr, pivotIndex + 1, right, k, \ldots)}
        \]
    \end{itemize}
\end{itemize}

\subsubsection{Complejidad}

La complejidad temporal promedio es:

\[
O(n)
\]

ya que en cada iteración solo se continúa explorando una partición del arreglo, y el tamaño del subarreglo se reduce aproximadamente por un factor constante en promedio.

Sin embargo, el peor caso ocurre cuando el pivote siempre queda en un extremo, resultando en:

\[
O(n^2)
\]

aunque este comportamiento es poco frecuente en datos aleatorios y mitiga considerablemente el uso de mediana de tres.

La complejidad espacial es $O(\log n)$ por las llamadas recursivas apiladas en el stack.

\section{Aplicación práctica}

\subsection{Ranking de los mejores N deportistas}
Para obtener el ranking de los \(N\) mejores deportistas según puntaje, se implementó la función \texttt{get\_top\_n\_athletes}, que recibe el arreglo de deportistas, su cantidad total y el valor \(N\). La función reserva un arreglo de \(N\) posiciones y lo completa llamando a \texttt{get\_kth\_athlete} con \(k = 1, 2, \ldots, N\), obteniendo cada puesto de forma independiente mediante Quick Select.

A diferencia de ordenar el arreglo completo, Quick Select localiza el \(k\)-ésimo elemento sin necesidad de ordenar todos los datos. Su costo promedio es:

\[
O(n)
\]

por llamada, lo que resulta en un costo total de:

\[
O(N \cdot n)
\]

para el ranking completo, frente al costo:

\[
O(n \log n)
\]

de un ordenamiento completo. Esto lo hace más eficiente cuando \(N\) es pequeño respecto al total de deportistas.

La función \texttt{show\_ranking} se encarga de cargar los datos desde el archivo CSV, invocar \texttt{get\_top\_n\_athletes} para obtener los resultados y mostrarlos utilizando las funciones de formato ya existentes en el sistema.
